\documentclass[11pt]{article}

\usepackage[a4paper,margin=1in]{geometry}
\usepackage{hyperref}
\usepackage{enumitem}
\usepackage{titlesec}

\titleformat{\section}{\large\bfseries}{\thesection}{1em}{}
\setlist[itemize]{noitemsep, topsep=4pt}
\setlist[enumerate]{noitemsep, topsep=4pt}

\title{KCPC-CS120 \\ Supplementary Reading \\ Combinatorics I: Learning to Count}
\author{}
\date{}

\begin{document}

\maketitle

\noindent
This document contains recommended and further readings for stars and bars, principle of inclusion-exclusion, bijections, and Burnside's Lemma (which requires knowledge of group theory). 

These references are intended to deepen conceptual understanding beyond implementation.

\hrule
\vspace{1em}

% ----------------------------------------------------
\section*{I. Core Algorithmic Reading}

\subsection*{Competitive Programmer's Handbook --- Antti Laaksonen (PDF)}
\url{https://cses.fi/book/book.pdf}
\vspace{0.5em}
\begin{itemize}
    \item \textbf{Counting} --- Chapter 22, p. 227
    \item \textbf{Binomial Coefficients} --- Section 22.2, p. 229
    \item \textbf{Inclusion-Exclusion} --- Section 22.4, p. 232
    \item \textbf{Burnside's Lemma} --- Section 22.6, p. 235
    \item \textbf{Catalan Numbers} --- Section 22.7, p. 236
\end{itemize}

\vspace{1em}
\hrule
\vspace{1em}

% ----------------------------------------------------
\section*{II. Online Reference --- CP-Algorithms}

\subsection*{Combinatorics}
\url{https://cp-algorithms.com/combinatorics/}
\vspace{0.5em}
\begin{itemize}
    \item \textbf{Binomial Coefficients} ---
    
        \url{https://cp-algorithms.com/combinatorics/binomial-coefficients.html}
    \item \textbf{Stars and Bars} ---
    
        \url{https://cp-algorithms.com/combinatorics/stars_and_bars.html}
    \item \textbf{Inclusion-Exclusion} ---
    
        \url{https://cp-algorithms.com/combinatorics/inclusion-exclusion.html}
    \item \textbf{Catalan Numbers} ---
    
        \url{https://cp-algorithms.com/combinatorics/catalan-numbers.html}
    \item \textbf{Burnside's Lemma} ---
    
        \url{https://cp-algorithms.com/combinatorics/burnside.html}
\end{itemize}

\vspace{1em}
\hrule
\vspace{1em}

% ----------------------------------------------------
\section*{III. Formal Textbook Reference}

\subsection*{Dexter Chua's Lecture Notes}
These are intended for a more advanced audience comfortable with proof-based mathematics or those willing to learn more about group theory or number theory.

Numbers and Sets:
Covers bijections, counting arguments, inclusion-exclusion, and
the formal foundations of combinatorics.

\url{hhttps://dec41.user.srcf.net/h/IA_M/numbers_and_sets}

Groups:
Required background for Burnside's Lemma. Read this only if you want
to understand the group-theoretic proof rather than just applying the formula.

\url{https://dec41.user.srcf.net/h/IA_M/groups}
\vspace{1em}

\hrule
\vspace{1em}

The ArtofProblemSolving has a lot of material on all sorts of olympiad content, recommended to check it out. It's intended more for people who've already done some olympiads though.

Brilliant is good for newbies and explains everything from the ground up.


\end{document}
